\begin{tikzpicture}[
    x=1pt,
    y=1pt,
    font=\small,
    >=stealth,
    panel/.style={draw=gray!35, rounded corners=14pt, fill=gray!4, line width=1pt},
    block/.style={draw=#1!75!black, fill=#1!12, rounded corners=9pt, line width=0.9pt, align=center},
    flow/.style={->, line width=0.95pt, draw=gray!75}
]
\def\W{742}
\def\H{270}
\fill[gray!4, rounded corners=16pt] (0,0) rectangle (\W,\H);
\draw[gray!35, rounded corners=16pt, line width=1pt] (0,0) rectangle (\W,\H);
\fill[cyan!18, rounded corners=14pt] (18,228) rectangle (724,256);
\node[anchor=west, font=\bfseries\small] at (34,242) {多轮越狱交互机制图};
\node[anchor=east, font=\scriptsize, text=gray!70!black] at (706,242) {在多轮语境中持续修正攻击策略};

\node[block=orange, minimum width=84pt, minimum height=32pt] (r1) at (96,152) {\textbf{第1轮输入}\\初始提示};
\node[block=blue, minimum width=84pt, minimum height=32pt] (m1) at (216,152) {\textbf{模型回复}\\拒答 / 暴露};
\node[block=teal, minimum width=84pt, minimum height=32pt] (j1) at (336,152) {\textbf{在线判定}\\当前状态};
\node[block=green, minimum width=84pt, minimum height=32pt] (u1) at (456,152) {\textbf{策略更新}\\生成追问};
\node[block=orange, minimum width=84pt, minimum height=32pt] (r2) at (576,152) {\textbf{第2轮输入}\\追加提示};
\node[block=blue, minimum width=84pt, minimum height=32pt] (m2) at (646,98) {\textbf{后续轮次}\\直至成功或终止};

\draw[flow] (r1.east) -- (m1.west);
\draw[flow] (m1.east) -- (j1.west);
\draw[flow] (j1.east) -- (u1.west);
\draw[flow] (u1.east) -- (r2.west);
\draw[flow] (r2.south east) .. controls (612,128) and (626,120) .. (m2.north west);
\draw[flow] (m2.south west) .. controls (590,64) and (486,58) .. node[below, font=\scriptsize] {继续迭代} (u1.south);

\draw[panel] (62,34) rectangle (682,72);
\node[anchor=west, font=\bfseries\scriptsize, text=red!70!black] at (78,58) {记录字段};
\node[anchor=west, font=\scriptsize] at (80,42) {每轮保存输入提示、模型回复、判定状态、策略信息, 并统计首次成功轮次与累计成功率};
\end{tikzpicture}
